\chapter{关于 Fourier 方法的推导}

前文中直接给出放大因子并据此检查数值格式的稳定性，显得非常突兀，这里尝试用线性代数以及一点泛函分析的语言进行推导，这种推导方式其实和从 Fourier 理论直接入手的做法殊途同归（参考张强《偏微分方程的有限差分方法》2.4 节）。

\section{周期边界情形}

考虑周期边界下的标量常系数单步格式
\[
    v_j^{n+1} = \sum_{\ell=-r}^{p} a_\ell v_{j+\ell}^n, \quad j=1,\cdots,N
\]
可以自然转换为向量形式
\[
    V^{n+1} = Q_N V^n = \cdots = (Q_N)^{n+1} V^0
\]
其中 $V^{n},V^{n+1} \in \mathbb{C}^N$ 是有限维向量，$Q_N \in \mathbb{C}^{N\times N}$ 满足
\[
    (Q_N V)_j = \sum_{\ell=-r}^{p} a_\ell\, V_{j+\ell},
\]
超过范围的下标在模 $N$ 的意义下理解。
定义循环移位矩阵
\[
    S_N =
    \begin{pmatrix}
        0 & 1 &        &        &   \\
          & 0 & 1      &        &   \\
          &   & \ddots & \ddots &   \\
          &   &        & 0      & 1 \\
        1 &   &        &        & 0
    \end{pmatrix} \in \mathbb{C}^{N\times N}
\]
那么 $Q_N$ 可以表示为
\[
    Q_N = \sum_{\ell=-r}^{p} a_\ell (S_N)^\ell,
\]

易得 $S_N$ 是酉方阵，可以酉相似对角化
\[
    S_N = W_N D_N W_N^*.
\]
其中 $W_N$ 是酉方阵，$D_N$ 是特征值组成的对角阵，具体为
\[
    D_N = \operatorname{diag}(e^{i\xi_1},\cdots,e^{i\xi_N}), \quad \xi_k=\frac{2\pi k}{N}
\]
因此 $Q_N$ 也可以酉相似对角化
\[
    Q_N = \sum_{\ell=-r}^{p} a_\ell (S_N)^\ell
    = W_N \left(\sum_{\ell=-r}^{p} a_\ell (D_N)^\ell\right) W_N^*
    = W_N \operatorname{diag}\bigl(q(e^{i\xi_1}),\dots,q(e^{i\xi_N})\bigr) W_N^*
\]
其中 Laurent 多项式
\[
    q(z) = \sum_{\ell=-r}^{p} a_\ell z^\ell.
\]
因此
\[
    \|(Q_N)^n\| = \max_{k=1,\dots,N} |q(e^{i\xi_k})|^n .
\]

数值格式的稳定性可以转换为如下要求：存在不依赖 $N$ 的常数 $C$，使得
\[
    \sup_{n\ge 0} \|(Q_N)^n\| \le C
\]
这里的记号 $\|\,\cdot\,\|$ 是 $L^2$ 范数诱导的算子范数。
对固定的 $N$，要求存在常数 $C_N$ 使得
\[
    \sup_{n\ge 0} \|(Q_N)^n\| \le C_N
\]
等价于要求
\[
    \max_{1\le k \le N} |q(e^{i\xi_k})| \le 1
\]
但是对于数值格式，我们关注的是网格加密后的极限情况，因此要求有界性对网格加密（$N\to\infty$）一致成立。
由 $q$ 的连续性，一致稳定性等价于
\[
    \sup_{\xi\in[-\pi,\pi]} |q(e^{i\xi})| \le 1.
\]
这就得到了严格的 von~Neumann 条件，并且 $q(e^{i\xi})$ 就是之前定义的放大因子。

\section{无界情形}

现在去掉周期边界，考虑整空间网格上的标量常系数单步格式
\[
    v_j^{n+1}
    =
    \sum_{\ell=-r}^{p}a_\ell v_{j+\ell}^n,
    \qquad j\in\mathbb Z.
\]
考虑无限维 Hilbert 空间
\[
    H_h=\ell_h^2(\mathbb Z)
    :=
    \left\{
        V=(V_j)_{j\in\mathbb Z}:
        h\sum_{j\in\mathbb Z}|V_j|^2<\infty
    \right\},
\]
其内积和范数分别为
\[
    \langle U,V\rangle_h
    =
    h\sum_{j\in\mathbb Z}U_j\overline{V_j},
    \qquad
    \|V\|_h
    =
    \left(
        h\sum_{j\in\mathbb Z}|V_j|^2
    \right)^{1/2}.
\]
定义双边移位算子
\[
    S:H_h\to H_h,
    \qquad
    (SV)_j=V_{j+1}.
\]
显然可逆并且逆算子为
\[
    (S^{-1}V)_j=V_{j-1}.
\]
定义演化算子
\[
    Q = q(S) = \sum_{\ell=-r}^{p}a_\ell S^\ell.
\]
因此数值格式可以表示为
\[
    V^{n+1} = Q V^n = \cdots = Q^{n+1} V^0
\]

双边移位算子 $S$ 是酉算子，谱为单位圆
\[
    \sigma(S)
    =
    \left\{
        e^{i \xi}:
        \xi\in[-\pi,\pi]
    \right\}
\]
由谱映射定理，$Q=q(S)$ 的谱为
\[
    \sigma(Q)
    =
    q(\sigma(S))
    =
    \left\{
        q(e^{i \xi}):
        \xi\in[-\pi,\pi]
    \right\}.
\]
由于 $S$ 是正规算子，所以 $Q=q(S)$ 也是正规算子，$Q^n$ 仍然正规，并且
\[
    \|Q^n\| = \max_{\lambda\in\sigma(Q^n)}|\lambda| = \max_{\lambda\in\sigma(Q)}|\lambda|^n = \max_{\xi\in[-\pi,\pi]}|q(e^{i\xi})|^n.
\]
重复之前的稳定性讨论，可以得到无界网格情形下的严格 von Neumann 条件
\[
    \sup_{\xi\in[-\pi,\pi]}
    |q(e^{i\xi})|
    \leq1.
\]
无界情况和周期边界的情况虽然有所区别，但是最终我们得到的是同样的结果。
